\documentclass[11pt,a4paper]{article}
\usepackage[utf8]{inputenc}
\usepackage[spanish,es-nodecimaldot,es-noquoting]{babel}
\usepackage{amsmath,amssymb,amsfonts}
\usepackage{geometry}
\geometry{top=1.8cm,bottom=1.8cm,left=1.8cm,right=1.8cm}
\usepackage{enumitem}
\usepackage{titlesec}
\usepackage{xcolor}
\usepackage{ragged2e}
\usepackage{booktabs}
\usepackage{tabularx}
\usepackage{tcolorbox}
\tcbuselibrary{skins,breakable}
\usepackage{hyperref}
\usepackage{fancyhdr}
\usepackage{listings}

% Paleta de colores institucional
\definecolor{unlbluelight}{RGB}{0,86,145}
\definecolor{azulOscuro}{RGB}{12,43,92}
\definecolor{verdeBosque}{RGB}{27,107,74}
\definecolor{rojoAlerta}{RGB}{179,38,30}
\definecolor{doradoUNL}{RGB}{197,143,39}
\definecolor{grisFondo}{RGB}{247,249,252}
\definecolor{grisBorde}{RGB}{209,217,230}
\definecolor{grisTerminal}{RGB}{30,30,30}

% Configuracion de hyperref
\hypersetup{
    colorlinks=true,
    linkcolor=azulOscuro,
    filecolor=magenta,      
    urlcolor=unlbluelight,
    citecolor=verdeBosque
}

% Configuracion de listings para terminal y codigo
\lstset{
  basicstyle=\footnotesize\ttfamily\color{black},
  backgroundcolor=\color{grisFondo},
  frame=single,
  rulecolor=\color{grisBorde},
  breaklines=true,
  breakatwhitespace=true,
  showstringspaces=false,
  tabsize=2,
  numbers=none,
  xleftmargin=6pt,
  xrightmargin=6pt,
  aboveskip=4pt,
  belowskip=4pt
}

% Columnas tabulares
\newcolumntype{Y}{>{\RaggedRight\arraybackslash}X}
\newcolumntype{C}[1]{>{\centering\arraybackslash}p{#1}}
\newcolumntype{L}[1]{>{\RaggedRight\arraybackslash}p{#1}}

% Cabeceras y pies de pagina
\setlength{\headheight}{22pt}
\pagestyle{fancy}
\fancyhf{}
\renewcommand{\headrulewidth}{0.8pt}
\renewcommand{\headrule}{\hbox to\headwidth{\color{unlbluelight}\leaders\hrule height \headrulewidth\hfill}}
\fancyhead[L]{\footnotesize\textbf{\color{unlbluelight}UNIVERSIDAD NACIONAL DE LOJA} $\vert$ Ing. Ambiental}
\fancyhead[R]{\footnotesize\textbf{Matemática II} $\vert$ Guía de Instalación y GitHub}
\fancyfoot[L]{\footnotesize Ciclo Sep 2026 -- Feb 2027}
\fancyfoot[C]{\footnotesize Guía Práctica de Terminal, Git, VS Code y Entorno Virtual}
\fancyfoot[R]{\footnotesize Pág. \thepage}

\fancypagestyle{plain}{%
  \fancyhf{}%
  \renewcommand{\headrulewidth}{0pt}%
  \fancyfoot[L]{\footnotesize Ciclo Sep 2026 -- Feb 2027}%
  \fancyfoot[C]{\footnotesize Guía Práctica de Terminal, Git, VS Code y Entorno Virtual}%
  \fancyfoot[R]{\footnotesize Pág. \thepage}%
}

% Formato de titulos
\titleformat{\section}{\large\bfseries\color{azulOscuro}}{\thesection.}{0.5em}{}[\titlerule]
\titleformat{\subsection}{\normalsize\bfseries\color{unlbluelight}}{\thesubsection.}{0.4em}{}
\titleformat{\subsubsection}{\small\bfseries\color{verdeBosque}}{\thesubsubsection.}{0.3em}{}
\titlespacing*{\section}{0pt}{9pt}{4pt}
\titlespacing*{\subsection}{0pt}{7pt}{3pt}
\titlespacing*{\subsubsection}{0pt}{5pt}{2pt}

% Cajas de diseno
\newtcolorbox{bannerbox}[1][]{
  colback=azulOscuro, colframe=azulOscuro, arc=4pt, boxrule=1pt,
  fonttitle=\bfseries\color{white}, colbacktitle=azulOscuro,
  top=6pt, bottom=6pt, left=8pt, right=8pt, breakable, #1
}

\newtcolorbox{conceptbox}[1][]{
  colback=grisFondo, colframe=unlbluelight, arc=3pt, boxrule=1pt,
  fonttitle=\bfseries\color{white}, colbacktitle=unlbluelight,
  top=5pt, bottom=5pt, left=8pt, right=8pt, breakable, #1
}

\newtcolorbox{stepbox}[2][]{
  colback=white, colframe=unlbluelight, arc=3pt, boxrule=1pt,
  title=\textbf{#2}, fonttitle=\bfseries\color{white}, colbacktitle=unlbluelight,
  top=5pt, bottom=5pt, left=8pt, right=8pt, breakable, #1
}

\newtcolorbox{alertbox}[1][]{
  colback=red!3!white, colframe=rojoAlerta, arc=3pt, boxrule=1pt,
  fonttitle=\bfseries\color{white}, colbacktitle=rojoAlerta,
  top=5pt, bottom=5pt, left=8pt, right=8pt, breakable, #1
}

\newtcolorbox{terminalbox}[2][]{
  colback=grisTerminal, colframe=black, arc=3pt, boxrule=1pt,
  title={\color{white}\footnotesize\textbf{Terminal / Consola: #2}},
  fonttitle=\bfseries\color{white}, colbacktitle=black,
  top=4pt, bottom=4pt, left=6pt, right=6pt, breakable, #1
}

\begin{document}
\thispagestyle{plain}

% ============================================================
% ENCABEZADO INSTITUCIONAL
% ============================================================
\begin{center}
  {\Large\textbf{\color{azulOscuro}UNIVERSIDAD NACIONAL DE LOJA}}\vspace{0.10cm}\\
  {\large\textbf{\color{unlbluelight}FACULTAD AGROPECUARIA Y DE RECURSOS NATURALES RENOVABLES}}\vspace{0.06cm}\\
  {\normalsize\textbf{CARRERA DE INGENIERÍA AMBIENTAL}}\vspace{0.10cm}\\
  {\small Período Académico: Septiembre 2026 -- Febrero 2027 $\quad\vert\quad$ Modalidad: Presencial}
\end{center}

\vspace{0.10cm}

\begin{bannerbox}
  \centering\color{white}
  {\Large\bfseries GUÍA PASO A PASO: REPOSITORIO GITHUB, VISUAL STUDIO CODE,}\\[0.10cm]
  {\Large\bfseries ENTORNO VIRTUAL PYTHON (venv) Y BASES DEL PROYECTO INTEGRADOR}\\[0.10cm]
  {\normalsize\textbf{MATEMÁTICA II: CÁLCULO INTEGRAL APLICADO A LA INGENIERÍA AMBIENTAL}}\\[0.06cm]
  {\footnotesize\textit{Docente: Dr. Guillermo Chuncho $\quad\vert\quad$ Repositorio: https://github.com/cguillermo79/Informacion\_matematica}}
\end{bannerbox}

\vspace{0.25cm}

% ============================================================
% 1. PROPOSITO DE LA GUIA
% ============================================================
\section{Propósito y Objetivos de la Guía}

Esta guía técnica tiene como finalidad orientar de manera detallada, secuencial y reproducible a los estudiantes de la Carrera de Ingeniería Ambiental en la configuración integral de su espacio de trabajo en sus computadoras personales. 

\begin{conceptbox}[title=Objetivos de Aprendizaje Operativo]
Al finalizar esta guía, cada estudiante y equipo de trabajo estará en capacidad de:
\begin{enumerate}[leftmargin=*,itemsep=1.5pt]
  \item Instalar y verificar las herramientas base de desarrollo científico: Git, Python 3 y Visual Studio Code.
  \item Clonar el repositorio oficial de GitHub del curso y mantenerlo actualizado localmente.
  \item Crear, configurar y activar un \textbf{entorno virtual aislado de Python (\texttt{venv})} para garantizar que los scripts y librerías funcionen sin colisiones de versiones.
  \item Vincular correctamente el intérprete virtual en Visual Studio Code.
  \item Acceder a las bases de datos geoespaciales y climáticas asignadas a su grupo (ventanas de estudio de $400$ o $100$ celdas en las zonas Norte/Sur y la serie cantonal de incendios del Cantón Loja).
  \item Conocer la estructura de carpetas donde desarrollarán y depositarán los entregables de la Unidad 1, Unidad 2 y Unidad 3, y el flujo de comandos en Git para sincronizar su trabajo.
\end{enumerate}
\end{conceptbox}

% ============================================================
% 2. REQUISITOS PREVIOS E INSTALACION DE HERRAMIENTAS BASE
% ============================================================
\section{Requisitos Previos: Instalación del Software Base}

Antes de ejecutar los comandos en la terminal, asegúrese de tener instalado en su computadora el siguiente software gratuito y de código abierto:

\begin{enumerate}[leftmargin=*,itemsep=4pt]
  \item \textbf{Git (Sistema de Control de Versiones):}
  \begin{itemize}
    \item Descargar desde el sitio oficial: \url{https://git-scm.com/downloads}.
    \item Durante la instalación en Windows, deje las opciones por defecto recomendadas (asegúrese de permitir Git desde la línea de comandos de Windows).
  \end{itemize}

  \item \textbf{Python 3 (Versión 3.10, 3.11 o 3.12 recomendada):}
  \begin{itemize}
    \item Descargar desde: \url{https://www.python.org/downloads/}.
    \item \textbf{\color{rojoAlerta}¡PASO CRÍTICO EN WINDOWS!:} Al iniciar el instalador, marque obligatoriamente la casilla inferior que dice: \textbf{\textit{''Add python.exe to PATH''}} antes de hacer clic en \textit{Install Now}. Si omite este paso, el comando \texttt{python} no será reconocido en la terminal.
  \end{itemize}

  \item \textbf{Visual Studio Code (VS Code):}
  \begin{itemize}
    \item Descargar desde: \url{https://code.visualstudio.com/Download}.
    \item Extensiones recomendadas dentro de VS Code (instalar desde la pestaña de Extensiones \texttt{Ctrl+Shift+X}):
    \begin{itemize}
      \item \textbf{Python} (desarrollada por Microsoft).
      \item \textbf{Jupyter} (para cuadernos interactivos \texttt{.ipynb}).
      \item \textbf{GitLens} (para trazabilidad y control de ramas en Git).
    \end{itemize}
  \end{itemize}
\end{enumerate}

% ============================================================
% 3. PASO A PASO EN LA TERMINAL: DE GITHUB A SU COMPUTADORA
% ============================================================
\section{Paso a Paso en la Terminal: Clonación y Apertura en VS Code}

Abra la terminal de su sistema operativo: en Windows puede usar \textbf{PowerShell} o la \textbf{Terminal de Windows}; en macOS o Linux utilice la aplicación \textbf{Terminal}.

\begin{stepbox}{Paso 1: Configurar su identidad en Git (Solo si es la primera vez)}
Abra su terminal y configure su nombre y su correo electrónico personal (el mismo asociado a su cuenta de GitHub) para que sus contribuciones queden registradas adecuadamente:
\begin{lstlisting}[language=bash]
git config --global user.name "Nombres y Apellidos del Estudiante"
git config --global user.email "correo.personal@gmail.com"
\end{lstlisting}
Para verificar que quedó bien configurado, ejecute:
\begin{lstlisting}[language=bash]
git config --list
\end{lstlisting}
\end{stepbox}

\begin{stepbox}{Paso 2: Crear una carpeta de trabajo y ubicarse en ella}
Se recomienda tener una carpeta ordenada en su disco local (por ejemplo en \texttt{Documentos}):
\begin{lstlisting}[language=bash]
# En Windows PowerShell o CMD:
cd $HOME\Documents
mkdir Matematica_UNL
cd Matematica_UNL

# En Linux o macOS:
cd ~/Documents
mkdir Matematica_UNL
cd Matematica_UNL
\end{lstlisting}
\end{stepbox}

\begin{stepbox}{Paso 3: Clonar el repositorio oficial del curso desde GitHub}
Ejecute el comando \texttt{git clone} utilizando la dirección URL pública del repositorio del curso:
\begin{lstlisting}[language=bash]
git clone https://github.com/cguillermo79/Informacion_matematica.git
\end{lstlisting}
Una vez que finalice la descarga de todos los archivos y bases de datos, ingrese a la carpeta del proyecto:
\begin{lstlisting}[language=bash]
cd Informacion_matematica
\end{lstlisting}
\end{stepbox}

\begin{stepbox}{Paso 4: Abrir el proyecto en Visual Studio Code}
Desde la misma terminal, dentro de la carpeta \texttt{Informacion\_matematica}, ejecute:
\begin{lstlisting}[language=bash]
code .
\end{lstlisting}
\textit{Nota:} El punto (\texttt{.}) indica que se debe abrir la carpeta actual. Si el comando \texttt{code} no abre VS Code, abra manualmente Visual Studio Code, vaya al menú superior: \textbf{File $\rightarrow$ Open Folder...} y seleccione la carpeta \texttt{Informacion\_matematica}.
\end{stepbox}

% ============================================================
% 4. CREACION Y ACTIVACION DEL ENTORNO VIRTUAL PYTHON (venv)
% ============================================================
\section{Creación y Activación del Entorno Virtual de Python (\texttt{venv})}

\begin{conceptbox}[title=¿Por qué es indispensable crear un entorno virtual?]
Un entorno virtual (\texttt{venv}) es una instalación aislada y autónoma de Python dentro del proyecto. Evita conflictos de paquetes entre diferentes asignaturas, garantiza que todas las librerías necesarias (\texttt{pandas}, \texttt{numpy}, \texttt{scipy}, \texttt{matplotlib}) tengan las versiones exactas y no modifica la configuración global de su sistema operativo.
\end{conceptbox}

Realice todos los siguientes pasos desde la \textbf{Terminal Integrada de Visual Studio Code} (puede abrirla pulsando las teclas \texttt{Ctrl + \~{}} o \texttt{Ctrl + Shift + \~{}}, o en el menú superior: \textbf{Terminal $\rightarrow$ New Terminal}):

\begin{stepbox}{Paso 5: Crear el entorno virtual en la raíz del proyecto}
Estando en la terminal dentro de \texttt{Informacion\_matematica}, ejecute:
\begin{lstlisting}[language=bash]
# En Windows (PowerShell o CMD):
python -m venv .venv

# En macOS o Linux (o si tiene multiples versiones instaladas):
python3 -m venv .venv
\end{lstlisting}
Esto creará una subcarpeta oculta llamada \texttt{.venv} en su proyecto con el intérprete de Python y su propio gestor de paquetes \texttt{pip}.
\end{stepbox}

\begin{alertbox}[title=Atención en Windows PowerShell: Error de Políticas de Ejecución]
Si al intentar activar el entorno virtual en Windows PowerShell le aparece el error:\\
\textit{\textcolor{rojoAlerta}{''File ... Activate.ps1 cannot be loaded because running scripts is disabled on this system''}},\\
debe habilitar la ejecución de scripts para su usuario ejecutando \textbf{una sola vez} este comando en la terminal:
\begin{lstlisting}[language=bash]
Set-ExecutionPolicy -ExecutionPolicy RemoteSigned -Scope CurrentUser
\end{lstlisting}
Presione \textbf{Y} (Sí) cuando el sistema le pida confirmación.
\end{alertbox}

\begin{stepbox}{Paso 6: Activar el entorno virtual}
Ejecute el comando correspondiente a su terminal y sistema operativo:
\begin{lstlisting}[language=bash]
# En Windows (PowerShell):
.\.venv\Scripts\Activate.ps1

# En Windows (CMD - Simbolo del sistema):
.\.venv\Scripts\activate.bat

# En macOS y Linux (Bash o Zsh):
source .venv/bin/activate
\end{lstlisting}
\textbf{¿Cómo saber si el entorno virtual está activo?} \\
En la terminal observará que al inicio de la línea de comandos aparece el prefijo \texttt{(.venv)} entre paréntesis, por ejemplo:
\begin{lstlisting}[language=bash]
(.venv) PS C:\Users\Estudiante\Documents\Matematica_UNL\Informacion_matematica>
\end{lstlisting}
\end{stepbox}

\begin{stepbox}{Paso 7: Seleccionar el Intérprete de Python dentro de Visual Studio Code}
Para que VS Code y los cuadernos de Jupyter reconozcan el entorno virtual:
\begin{enumerate}[leftmargin=*,itemsep=2pt]
  \item Presione en su teclado la combinación \texttt{Ctrl + Shift + P} (o \texttt{Cmd + Shift + P} en Mac).
  \item Escriba: \textbf{\texttt{Python: Select Interpreter}} y presione Enter.
  \item En la lista desplegable, elija la opción recomendada que contiene \texttt{(.venv): .venv\textbackslash Scripts\textbackslash python.exe} (o \texttt{.venv/bin/python}).
\end{enumerate}
En la esquina inferior derecha de la ventana de VS Code verá reflejado \texttt{3.1x.x ('.venv': venv)}.
\end{stepbox}

\begin{stepbox}{Paso 8: Actualizar pip e instalar las librerías científicas requeridas}
Con el entorno virtual activo \texttt{(.venv)}, actualice el gestor \texttt{pip} e instale todas las librerías científicas del curso ejecutando:
\begin{lstlisting}[language=bash]
python -m pip install --upgrade pip
pip install -r requirements.txt
\end{lstlisting}
\textit{Si no tuviera el archivo \texttt{requirements.txt}, puede instalar las librerías directamente con:}
\begin{lstlisting}[language=bash]
pip install numpy scipy pandas matplotlib seaborn jupyter pyarrow fastparquet openpyxl
\end{lstlisting}
\end{stepbox}

% ============================================================
% 5. BASES DE DATOS DEL PROYECTO INTEGRADOR POR GRUPO
% ============================================================
\section{Acceso y Verificación de las Bases de Datos por Equipo}

Las bases de datos oficiales ya se encuentran descargadas en la carpeta \texttt{Unidad\_1/\allowbreak bases\_\allowbreak datos\_\allowbreak proyecto\_\allowbreak integrador/} (archivos CSV pequeños, de 5\,KB a 1\,MB). El significado de cada columna está en \texttt{DICCIONARIO\_\allowbreak VARIABLES.csv}. Cada equipo debe identificar la subcarpeta asignada a su temática y zona de estudio:

\begin{table}[h!]
\centering
\scriptsize
\renewcommand{\arraystretch}{1.15}
\begin{tabularx}{\textwidth}{C{1.2cm} L{3.5cm} C{2.0cm} Y L{4.5cm}}
\toprule
\textbf{Equipo} & \textbf{Bloque Temático Asignado} & \textbf{Zona} & \textbf{Carpeta Asignada} & \textbf{Dimensión de la Base} \\
\midrule
\textbf{Equipo 1} & Relieve y Accesibilidad & Norte & \texttt{equipo\_01\_relieve\_}\newline\texttt{norte/} & 400 celdas (20$\times$20), 400 reg, 17 col \\
\textbf{Equipo 2} & Relieve y Accesibilidad & Sur & \texttt{equipo\_02\_relieve\_}\newline\texttt{sur/} & 400 celdas (20$\times$20), 400 reg, 17 col \\
\textbf{Equipo 3} & Cobertura Vegetal y Suelo & Norte & \texttt{equipo\_03\_cobertura\_}\newline\texttt{norte/} & 100 celdas (10$\times$10) $\times$ 72 meses, 7.200 reg, 28 col \\
\textbf{Equipo 4} & Cobertura Vegetal y Suelo & Sur & \texttt{equipo\_04\_cobertura\_}\newline\texttt{sur/} & 100 celdas (10$\times$10) $\times$ 72 meses, 7.200 reg, 28 col \\
\textbf{Equipo 5} & Clima y Atmósfera & Norte & \texttt{equipo\_05\_clima\_}\newline\texttt{norte/} & 100 celdas (10$\times$10) $\times$ 72 meses, 7.200 reg, 22 col \\
\textbf{Equipo 6} & Clima y Atmósfera & Sur & \texttt{equipo\_06\_clima\_}\newline\texttt{sur/} & 100 celdas (10$\times$10) $\times$ 72 meses, 7.200 reg, 22 col \\
\textbf{Equipo 7} & Incendios y Severidad & Cantonal & \texttt{equipo\_07\_incendios\_}\newline\texttt{loja/} & Serie mensual (72 reg) y resumen por celda (7.997 reg) \\
\bottomrule
\end{tabularx}
\caption{Subcarpetas autorizadas para cada equipo de trabajo en el proyecto integrador.}
\label{tab:equipos_carpetas}
\end{table}

\begin{stepbox}{Paso 9: Script de verificación rápida en Python}
Para verificar que el entorno virtual lee correctamente la base de datos de su equipo, cree un archivo de prueba \texttt{test\_datos.py} o abra un cuaderno de Jupyter y ejecute el siguiente código (modifique la ruta según su equipo):
\begin{lstlisting}[language=python]
import pandas as pd
from pathlib import Path

# Ejemplo para el Equipo 1 (Relieve Zona Norte):
carpeta = Path("Unidad_1/bases_datos_proyecto_integrador/equipo_01_relieve_norte")
ruta_base = carpeta / "equipo_01_relieve_norte.csv"

if ruta_base.exists():
    df = pd.read_csv(ruta_base)
    print("Lectura exitosa de la base de datos!")
    print(f"Filas totales: {df.shape[0]}, Columnas: {df.shape[1]}")
    print(f"Celdas espaciales unicas: {df['cell_id'].nunique()}")
    print(df.head(3))
else:
    print(f"Error: No se encontro el archivo en {ruta_base}")
\end{lstlisting}
\end{stepbox}

% ============================================================
% 6. ESTRUCTURA DE CARPETAS PARA ENTREGABLES Y SUBIDAS
% ============================================================
\section{Estructura de Carpetas de Entregables por Unidad}

El repositorio cuenta con una estructura jerárquica estandarizada para las tres unidades del ciclo académico. Cada equipo debe colocar sus trabajos exclusivamente en las carpetas correspondientes:

\begin{lstlisting}[language=bash]
Informacion_matematica/
|-- Unidad_1/
|   |-- bases_datos_proyecto_integrador/ <- Bases oficiales (solo lectura)
|   |-- 01_formulacion_3pct/             <- Fase 1: Delimitacion, pregunta, hipotesis y variables
|   |-- 02_avance_4pct/                  <- Fase 2: Scripts, cuadernos exploratorios y tablas
|   |-- 03_producto_tecnico_10pct/       <- Fase 3: Informe formal escrito en LaTeX y PDF
|   |-- 04_defensa_individual_8pct/      <- Fase 4: Presentacion y diapositivas de sustentacion
|   `-- 05_Guias/                        <- Guias oficiales de catedra (.pdf y .tex)
|-- Unidad_2/
|   |-- 01_formulacion_3pct/             <- Fase 1 U2: Modelos continuos no lineales
|   |-- 02_avance_4pct/                  <- Fase 2 U2: Tecnicas de integracion computacionales
|   |-- 03_producto_tecnico_10pct/       <- Fase 3 U2: Informe de Metodologia en LaTeX y PDF
|   |-- 04_defensa_individual_8pct/      <- Fase 4 U2: Material de sustentacion
|   `-- 05_Guias/
`-- Unidad_3/
    |-- 01_formulacion_3pct/             <- Fase 1 U3: Modelacion fisica y balances
    |-- 02_avance_4pct/                  <- Fase 2 U3: Integrales dobles y serie temporal
    |-- 03_producto_tecnico_10pct/       <- Fase 3 U3: Informe Final Consolidado (LaTeX y PDF)
    |-- 04_defensa_individual_8pct/      <- Fase 4 U3: Defensa final individual
    `-- 05_Guias/
\end{lstlisting}

% ============================================================
% 7. FLUJO DE TRABAJO EN GIT: ACTUALIZAR Y SUBIR TRABAJOS
% ============================================================
\section{Flujo de Trabajo con Git: Actualizar y Subir Entregables}

Para sincronizar el trabajo de su equipo y mantener al día las guías y materiales compartidos por el docente:

\begin{stepbox}{1. Descargar actualizaciones del docente (Antes de empezar a trabajar)}
Cada vez que el docente publique nuevas guías o avisos en el repositorio oficial, ejecute en su terminal:
\begin{lstlisting}[language=bash]
git pull origin main
\end{lstlisting}
Esto incorporará los nuevos archivos a su computador sin alterar su código local.
\end{stepbox}

\begin{stepbox}{2. Registrar y guardar los avances de su equipo}
Una vez que hayan elaborado su documento de formulación, script o informe en la carpeta correspondiente de su unidad:
\begin{lstlisting}[language=bash]
# Verifique que archivos modifico o creo:
git status

# Agregue los archivos de su entrega (ejemplo para Unidad 1):
git add Unidad_1/01_formulacion_3pct/

# Guarde el commit con un mensaje claro y formal:
git commit -m "Equipo 01: Entrega Fase 1 Formulacion del Proyecto Unidad 1"
\end{lstlisting}
\end{stepbox}

\begin{stepbox}{3. Subir sus cambios al repositorio remoto}
Para enviar sus cambios a GitHub (si trabaja en una rama de equipo o en su fork):
\begin{lstlisting}[language=bash]
git push origin main
\end{lstlisting}
\end{stepbox}

\begin{alertbox}[title=Buenas Prácticas Obligatorias: Archivo .gitignore]
Nunca suba a Git archivos pesados autogenerados ni entornos virtuales. El repositorio cuenta con un archivo \texttt{.gitignore} configurado para omitir automáticamente:
\begin{itemize}[leftmargin=*,itemsep=1pt]
  \item La carpeta del entorno virtual: \texttt{.venv/}
  \item Cachés de Python: \texttt{\_\_pycache\_\_/}, \texttt{*.pyc}
  \item Puntos de control de Jupyter: \texttt{.ipynb\_checkpoints/}
  \item Archivos temporales de compilación de LaTeX: \texttt{*.aux}, \texttt{*.log}, \texttt{*.out}, \texttt{*.toc}, \texttt{*.synctex.gz} (solo se suben el \texttt{.tex} y el \texttt{.pdf}).
\end{itemize}
\end{alertbox}

% ============================================================
% 8. RESUMEN DE COMANDOS MAS FRECUENTES (CHEATSHEET)
% ============================================================
\section{Resumen Rápido de Comandos en Terminal (Cheatsheet)}

\begin{table}[h!]
\centering
\small
\renewcommand{\arraystretch}{1.20}
\begin{tabularx}{\textwidth}{L{4.2cm} L{5.5cm} Y}
\toprule
\textbf{Acción Requerida} & \textbf{Comando Windows (PowerShell)} & \textbf{Comando Linux / macOS} \\
\midrule
Clonar repositorio & \texttt{git clone <URL>} & \texttt{git clone <URL>} \\
Abrir carpeta en VS Code & \texttt{code .} & \texttt{code .} \\
Crear entorno virtual & \texttt{python -m venv .venv} & \texttt{python3 -m venv .venv} \\
Habilitar scripts (solo 1 vez) & \texttt{Set-ExecutionPolicy RemoteSigned -Scope CurrentUser} & No aplica \\
Activar entorno virtual & \texttt{.\textbackslash .venv\textbackslash Scripts\textbackslash Activate.ps1} & \texttt{source .venv/bin/activate} \\
Desactivar entorno virtual & \texttt{deactivate} & \texttt{deactivate} \\
Instalar librerías del curso & \texttt{pip install -r requirements.txt} & \texttt{pip install -r requirements.txt} \\
Actualizar cambios remotos & \texttt{git pull origin main} & \texttt{git pull origin main} \\
Ver estado de cambios & \texttt{git status} & \texttt{git status} \\
Preparar archivos & \texttt{git add <carpeta o archivo>} & \texttt{git add <carpeta o archivo>} \\
Confirmar entrega local & \texttt{git commit -m "Mensaje"} & \texttt{git commit -m "Mensaje"} \\
Subir cambios al remoto & \texttt{git push origin main} & \texttt{git push origin main} \\
\bottomrule
\end{tabularx}
\caption{Referencia rápida de comandos de terminal para el proyecto integrador.}
\label{tab:cheatsheet}
\end{table}

\vfill
\begin{center}
  \footnotesize\textbf{Universidad Nacional de Loja $\vert$ Facultad Agropecuaria y de Recursos Naturales Renovables}\\
  \footnotesize Carrera de Ingeniería Ambiental $\vert$ Matemática II (Cálculo Integral) $\vert$ Ciclo Sep 2026 -- Feb 2027
\end{center}

\end{document}
